\section{Congruences and constructive periods}
\label{sec:congruences}

All congruences in this section are decidable integer assertions.
Polynomial calculations modulo $m$ mean coefficient-list calculations with
entries in $\{0,\ldots,m-1\}$. For a polynomial $F(X)=\sum_jc_jX^j$, use the
finite shorthand
\[
 (F(E)B)_n=\sum_jc_jB_{n+j}.
\]
The letter $E$ denotes a shift in the subscript, not an additional object
quantified over in the arithmetic language.

\subsection{Touchard's congruence}

\begin{theorem}[Touchard]\label{thm:touchard}
If $p$ is prime and $n\geq0$, then
\begin{equation}\label{eq:touchard}
 B_{n+p}\equiv B_{n+1}+B_n\pmod p.
\end{equation}
In particular $B_p\equiv2\pmod p$.
\end{theorem}
\begin{proof}
Let a cyclic permutation rotate $p$ distinguished elements of an
$(n+p)$-element set and fix the other $n$ elements. It acts, by a bounded
relabel-and-normalize algorithm, on the finite codes of partitions. Each
orbit has size one or $p$: the least positive return time divides the prime
$p$. One can find that return time by testing the $p$ possibilities. Choosing
the least code of each orbit therefore shows that, modulo $p$, only fixed
partitions contribute to their total number.

We classify the fixed partitions. A block meeting a fixed element must be
fixed as a block, since its rotate contains the same fixed element. If that
block meets the rotating set, it consequently contains the entire rotating
set. More generally, the restriction of a fixed partition to the $p$
rotating elements is a translation-invariant equivalence relation on
$\mathbb Z/p\mathbb Z$. The class of zero is closed under differences:
if $a$ and $b$ are equivalent to zero, translating $a\sim b$ by $-b$ gives
$a-b\sim0$. Hence it is a subgroup. A subgroup containing a nonzero element
contains all residues, because its first $p$ multiples are distinct;
otherwise it is $\{0\}$. Thus the rotating elements are either all in one
block or all in singleton blocks.

In the first case, collapse their common block contribution to one
distinguished element; this gives exactly $B_{n+1}$ partitions, allowing that
element to join any fixed elements. In the singleton case no rotating element
can share a block with a fixed element, by the preceding observation, leaving
exactly $B_n$ choices. The classes are disjoint. This proves
\eqref{eq:touchard}; setting $n=0$ gives the last assertion.
\end{proof}

There is also a short algebraic proof that is particularly convenient for
formalization. Lemma~\ref{lem:finite-shift} gives
\[
 \mathcal L\bigl(((X)_p-1)X^n\bigr)
  =\mathcal L((X+p)^n-X^n)\equiv0\pmod p.
\]
In $\mathbb F_p[X]$ one has
\begin{equation}\label{eq:finite-field-product}
 (X)_p=\prod_{a=0}^{p-1}(X-a)=X^p-X.
\end{equation}
To verify this last equality, the binomial theorem proves
$(a+1)^p\equiv a^p+1\pmod p$; induction through the residues gives
$a^p=a$ for all $a\in\mathbb F_p$. Both monic polynomials in
\eqref{eq:finite-field-product} have degree $p$ and these $p$ roots. Repeated
division by $X-a$ is legitimate because distinct residues have nonzero,
invertible differences, so their difference of degree less than $p$ is zero.
This proves \eqref{eq:touchard} by finite coefficient manipulation.

\subsection{Prime-power shifts and the prime period}

\begin{theorem}[Generalized Touchard congruence]\label{thm:touchard-power}
For prime $p$ and $r,n\geq0$,
\begin{equation}\label{eq:touchard-power}
 B_{n+p^r}\equiv B_{n+1}+rB_n\pmod p.
\end{equation}
\end{theorem}
\begin{proof}
Work with remainder polynomials modulo $Q(X)=X^p-X-1$ over
$\mathbb F_p$. Each remainder has degree less than $p$ and is a finite list
of residues. The binomial theorem in characteristic $p$ says
$(U+V)^p=U^p+V^p$. Induction on $r$ therefore gives
\[
 X^{p^r}=X+r\quad\text{in }\mathbb F_p[X]/(Q):
\]
the case $r=0$ is immediate; raising to the $p$th power replaces $X+r$ by
$X^p+r=X+(r+1)$, using the elementary residue identity $r^p=r$.
Thus $X^{p^r}-X-r$ is a multiple of $Q$. Every shift of
\eqref{eq:touchard} says $(Q(E)B)_n=0$ modulo $p$; multiplying by a polynomial
and taking the corresponding finite linear combination of these shifts
preserves zero. This gives \eqref{eq:touchard-power}.
\end{proof}

\begin{theorem}[An explicit period modulo a prime]\label{thm:prime-period}
For every prime $p$, let
\[
  N_p=1+p+\cdots+p^{p-1}=\frac{p^p-1}{p-1}.
\]
Then, for all $n\geq0$,
\begin{equation}\label{eq:prime-period}
 B_{n+N_p}\equiv B_n\pmod p.
\end{equation}
The least positive period exists by bounded search and divides $N_p$.
\end{theorem}
\begin{proof}
In the same finite quotient ring,
\[
 X^{N_p}=\prod_{r=0}^{p-1}X^{p^r}
       =\prod_{r=0}^{p-1}(X+r)=X^p-X=1.
\]
The middle product equals $X^p-X$ by
\eqref{eq:finite-field-product}, since $r\mapsto-r$ permutes the residues.
Hence $Q$ divides $X^{N_p}-1$, and the shift calculation just used proves
\eqref{eq:prime-period}. No field-extension existence theorem or
irreducibility theorem is required: the ring consists explicitly of $p^p$
remainder lists and need not first be shown to be a field.

A candidate period $t$ can be tested by the $p$ congruences
$B_{t+i}\equiv B_i$ for $0\leq i<p$. The order-$p$ recurrence
\eqref{eq:touchard} propagates those congruences to every subsequent index.
Search $1\leq t\leq N_p$ for the least passing value $T$; the preceding
paragraph supplies a passing value. Divide $N_p=qT+s$, $0\leq s<T$.
Repeated use of period $T$ and period $N_p$ gives
$B_{n+s}\equiv B_n$ for every $n$. If $s>0$, it passes the same finite test,
contrary to minimality. Therefore $s=0$.
\end{proof}

For $p=2$, $N_p=3$ and the initial residues are $1,1,0$. Thus
\begin{equation}\label{eq:bell-parity}
 B_n\text{ is even}\quad\Longleftrightarrow\quad n\equiv2\pmod3.
\end{equation}
For $p=3$, the analogous finite test shows that the least period is $13$:
$13$ is prime, a period divides it, and period one fails since
$B_1\not\equiv B_2\pmod3$. The statement that $N_p$ is the least period for
every prime is not proved here; the proof establishes that it is a period
and that the least period divides it.

\subsection{Periodicity for every modulus}

\begin{theorem}[A primitive recursive period for every modulus]
\label{thm:all-moduli-period}
For every integer $m\geq1$, the Bell sequence is purely periodic modulo $m$.
There is a period $T$ with $1\leq T\leq m^m$, found by a bounded computation.
In particular the following uniform period is valid:
\begin{equation}\label{eq:all-moduli-period}
 B_{n+(m^m)!}\equiv B_n\pmod m\qquad(n\geq0).
\end{equation}
\end{theorem}
\begin{proof}
For every polynomial $f\in\mathbb Z[X]$, the binomial theorem gives
$f(X+m)\equiv f(X)\pmod m$, coefficientwise. Hence the finite shift identity
gives
\[
 \mathcal L\bigl(((X)_m-1)f(X)\bigr)
   =\mathcal L(f(X+m)-f(X))\equiv0\pmod m.
\]
Here $\mathcal L$ takes integer-coefficient polynomials to integers, because
its values on monomials are the integer Bell numbers. Write
\[
 Q_m(X)=(X)_m-1=X^m+c_{m-1}X^{m-1}+\cdots+c_1X-1.
\]
Taking $f=X^n$ yields the order-$m$ recurrence
\begin{equation}\label{eq:composite-recurrence}
 B_{n+m}\equiv B_n-\sum_{j=1}^{m-1}c_jB_{n+j}\pmod m.
\end{equation}
Its state vector is
$v_n=(B_n,\ldots,B_{n+m-1})\bmod m$, an element of a list of $m^m$ states.
The transition shifts the entries left and appends the right side of
\eqref{eq:composite-recurrence}. It is a bijection: from the new state,
recover the discarded entry by adding
$\sum_{j=1}^{m-1}c_jB_{n+j}$ to its last entry.

Among $v_0,\ldots,v_{m^m}$ two states agree, by the finite pigeonhole
principle. That principle here is proved by scanning the codes
$0,\ldots,m^m-1$: if no code repeats, the number of entries cannot exceed the
number of codes. If $v_a=v_b$ with $0\leq a<b\leq m^m$, apply the explicit
inverse transition $a$ times to get $v_0=v_{b-a}$. Therefore
$T=b-a$ is a period from the start; the forward recurrence proves
$v_{n+T}=v_n$ for every $n$. Such a $T$ is found by scanning
$1,\ldots,m^m$ for a return to $v_0$. Since every positive integer at most
$m^m$ divides $(m^m)!$, repeating the period proves
\eqref{eq:all-moduli-period}. The case $m=1$ is included, with its unique
residue state.
\end{proof}

The general bound is deliberately generous. Its purpose is to give a
uniform primitive recursive construction, while
Theorem~\ref{thm:prime-period} provides the much sharper bound at prime
moduli. The accompanying exact-arithmetic program certifies the source's
least periods for moduli $1,\ldots,12$:
\[
 1,3,13,12,781,39,137257,24,39,2343,28531167061,156.
\]
For every listed modulus other than $11$ it checks the complete orbit up to
its first return. For $11$ it uses the shorter polynomial certificate
\[
 N_{11}=28531167061=15797\cdot1806113,
 \qquad B_{15797}\equiv6,\quad B_{1806113}\equiv9\pmod{11}.
\]
Bounded trial division verifies primality of the two factors. Binary powering
with the reduction rule $X^{11}=X+1$ verifies the residues; the output includes
both remainder coefficient lists. Since $B_0\equiv1$, neither maximal proper
divisor is a period. The least period divides $N_{11}$, so it must equal
$N_{11}$. These are finite computation certificates with a specified checker,
not claims that the accompanying Python program is an \ipra\ proof checker.
Checks at selected moduli do not prove a universal least-period conjecture.

\paragraph{Arithmetic audit.}
Primality is the bounded test that no $2\leq d<p$ divides $p$. The proof of
Touchard's congruence uses at most $p$ iterates per finite partition code.
The quotient-ring proof uses lists of length $p$ with residue coefficients;
raising to $p^r$ can be replaced by $r$ successive Frobenius steps.
The general period proof scans at most $m^m+1$ state codes, each with $m$
entries. Existence of a period therefore comes with a primitive recursive
witness. The propagation arguments are quantifier-free inductions on the
decidable congruences or on their bounded conjunctions.
